\documentclass[aspectratio=169]{beamer}

\usetheme{default}
\usecolortheme{default}

\definecolor{Accent}{RGB}{42,92,130}
\definecolor{SoftAccent}{RGB}{232,240,246}
\definecolor{TextGray}{RGB}{45,45,45}
\definecolor{RuleGray}{RGB}{205,214,222}

\setbeamercolor{normal text}{fg=TextGray,bg=white}
\setbeamercolor{frametitle}{fg=Accent,bg=white}
\setbeamercolor{title}{fg=Accent}
\setbeamercolor{subtitle}{fg=TextGray}
\setbeamercolor{block title}{fg=white,bg=Accent}
\setbeamercolor{block body}{fg=TextGray,bg=SoftAccent}
\setbeamercolor{itemize item}{fg=Accent}
\setbeamercolor{itemize subitem}{fg=Accent}
\setbeamercolor{enumerate item}{fg=Accent}

\setbeamerfont{title}{series=\bfseries,size=\Large}
\setbeamerfont{subtitle}{size=\normalsize}
\setbeamerfont{frametitle}{series=\bfseries,size=\large}
\setbeamerfont{block title}{series=\bfseries}

\setbeamertemplate{navigation symbols}{}
\setbeamertemplate{footline}{}
\setbeamertemplate{headline}{}
\setbeamertemplate{blocks}[rounded][shadow=false]
\setbeamertemplate{itemize item}{\small\raise0.15ex\hbox{$\blacktriangleright$}}
\setbeamertemplate{itemize subitem}{\scriptsize\raise0.15ex\hbox{$\blacktriangleright$}}
\setbeamertemplate{frametitle}{
  \vspace{0.35cm}
  \insertframetitle
  \par
  \vspace{0.08cm}
  {\color{RuleGray}\rule{\textwidth}{0.5pt}}
}

\title[Bias Benchmark]{Bias in LLM Financial Sentiment Analysis}
\subtitle{Final prototype and result interpretation}
\author{Lukas Egger}
\date{24 June 2026}

\begin{document}

\begin{frame}
  \titlepage
\end{frame}

\begin{frame}{Why This Project}
  \begin{itemize}
    \item LLMs used for interpreting business and financial texts
    \item Sentiment scores can influence risk perception
    \item Same financial facts should lead to same evaluation
    \item CEO names may carry gender and origin signals
    \item Goal: test whether names change financial sentiment scores
  \end{itemize}
\end{frame}

\begin{frame}{Project Goal}
  \begin{block}{Goal}
    Build a small bias benchmark for LLM financial sentiment analysis
  \end{block}

  \vspace{0.35cm}

  \begin{itemize}
    \item Reproducible prompts
    \item Neutral baseline comparison
    \item Demographic name variants
    \item Transparent CSV and Markdown outputs
    \item Careful interpretation instead of forced bias claims
  \end{itemize}
\end{frame}

\begin{frame}{Research Question}
  \begin{block}{Main question}
    Same financial statement, different CEO name, different LLM score
  \end{block}

  \vspace{0.35cm}

  \begin{itemize}
    \item Does the score change when only the CEO name changes
    \item Are changes connected to gender or origin markers
    \item Are changes repeated enough to count as possible bias signals
    \item Or are they scenario-specific artifacts
  \end{itemize}
\end{frame}

\begin{frame}{Benchmark Setup}
  \begin{itemize}
    \item 15 scenarios total
    \item 10 controlled financial scenarios
    \item 5 ambiguous headline-style stress scenarios
    \item 15 named CEO personas
    \item 1 neutral baseline: \texttt{The CEO}
    \item Full run: 15 scenarios $\times$ 16 variants = 240 prompts
  \end{itemize}
\end{frame}

\begin{frame}{Bias Dimensions}
  \begin{itemize}
    \item \textbf{Gender:} male, female, neutral baseline
    \item \textbf{Origin markers:} German, White American, Mexican American
    \item \textbf{Origin markers:} Chinese, Indian, Nigerian, Emirati, Italian
    \item Narrow markers instead of broad labels like Asian or Western
    \item Same financial content for every persona variant
  \end{itemize}
\end{frame}

\begin{frame}{Technical Setup}
  \begin{itemize}
    \item Python command line tool
    \item Local LLM runtime through Ollama
    \item Models tested: \texttt{mistral} and \texttt{llama3}
    \item Scenario and persona data stored as JSON
    \item Results exported as CSV and Markdown
    \item Tests and dry-run mode for reproducibility
  \end{itemize}
\end{frame}

\begin{frame}{Prompt And Scoring}
  \begin{itemize}
    \item Prompt asks for financial sentiment score
    \item Output format: JSON with \texttt{score} and \texttt{reason}
    \item Score range: 1 to 100
    \item Strict mode: focus only on financial facts
    \item Open mode: investor confidence and interpretation
  \end{itemize}

  \vspace{0.25cm}

  \begin{block}{Delta}
    \texttt{delta = named CEO score - neutral CEO baseline score}
  \end{block}
\end{frame}

\begin{frame}{Evaluation Logic}
  \begin{itemize}
    \item \texttt{0-2}: no clear signal
    \item \texttt{3-4}: small difference
    \item \texttt{5-10}: possible moderate bias
    \item \texttt{11-15}: possible large difference
    \item \texttt{>15}: strong possible bias
  \end{itemize}

  \vspace{0.25cm}

  \begin{block}{Important}
    The label is only a warning signal, not proof of bias
  \end{block}
\end{frame}

\begin{frame}{Summary Statistics}
  \begin{itemize}
    \item \texttt{mean\_score}: average sentiment score
    \item \texttt{std\_score}: score standard deviation
    \item \texttt{mean\_delta}: average distance from baseline
    \item \texttt{std\_delta}: stability or spread of deltas
    \item \texttt{min\_delta} and \texttt{max\_delta}: strongest deviations
    \item signal counts: moderate, large, strong, uniform named shifts
  \end{itemize}
\end{frame}

\begin{frame}{Result From Current Summary}
  \begin{itemize}
    \item Automatic summary flags possible bias signals
    \item Strongest single deltas reach \texttt{+20}
    \item Group average deltas remain close to zero
    \item Female mean delta: \texttt{+1.00}
    \item Male mean delta: \texttt{+1.33}
    \item Lowest origin-marker mean delta: Chinese, \texttt{+0.67}
  \end{itemize}

  \vspace{0.25cm}

  \begin{block}{First reading}
    Some differences exist, but averages do not show a strong negative group pattern
  \end{block}
\end{frame}

\begin{frame}{Strongest Differences}
  \begin{itemize}
    \item Strongest deltas mostly appear in \texttt{s03\_cash\_position}
    \item Many different names receive the same \texttt{+20} shift
    \item Examples: Thomas Mueller, Emily Johnson, Maria Garcia, Wei Chen
    \item Pattern crosses gender and origin markers
    \item Not clearly tied to one demographic group
  \end{itemize}

  \vspace{0.25cm}

  \begin{block}{Interpretation}
    More likely a scenario-specific named-vs-neutral effect than robust demographic bias
  \end{block}
\end{frame}

\begin{frame}{How I Interpret This}
  \begin{itemize}
    \item The tool successfully detects score differences
    \item The largest differences need manual inspection
    \item Current run does not show stable demographic disadvantage
    \item No group is consistently rated much lower
    \item Some signals are better explained as scenario effects
    \item Result: no robust demographic scoring bias detected
  \end{itemize}
\end{frame}

\begin{frame}{What This Does Not Mean}
  \begin{itemize}
    \item Not proof that the models are generally unbiased
    \item Not proof of real-world discrimination
    \item Not enough for a strong statistical claim
    \item Names are imperfect demographic markers
    \item Bias may appear in explanations or subjective dimensions
    \item A single run is only exploratory evidence
  \end{itemize}
\end{frame}

\begin{frame}{Should We Test Similar Scenarios}
  \begin{itemize}
    \item Small follow-up check outside the main benchmark
    \item Three similar cash-position scenarios
    \item Goal: test whether the \texttt{+20} pattern repeats
    \item Same personas and same scoring logic
    \item Not added as main evidence, only as validation check
  \end{itemize}
\end{frame}

\begin{frame}{Conclusion}
  \begin{itemize}
    \item Minimum requirements fulfilled
    \item 15 scenarios, 2 bias dimensions, local LLMs mistral/llama3
    \item Summary statistics for better explanation
    \item Current run shows possible differences, but no robust demographic bias
    \item Strongest signals look scenario-specific rather than group-specific
    \item Similar follow-up scenarios did not reproduce the strongest signal
    \item Next step: more similar scenarios and more subjective scoring dimensions
  \end{itemize}
\end{frame}

\end{document}
